\section{Configurational Entropy and the Equilibrium Concentration}

% ------------------------------------------------------------------
% Duy's derivation with every intermediate step kept.
% Lines marked "% CHANGED:" correct an error in the draft; the comment
% says what changed. Everything else follows the draft.
% ------------------------------------------------------------------

\paragraph{Counting the arrangements.}
A crystal has $N$ sites and $n$ defects to occupy them. Each defect has $g_X$
internal orientations. The problem is to count the ways to put $n$ defects on
$N$ sites, with $g_X$ orientations for each defect.
% CHANGED: the draft counted n defects on g_X N "orientational sites", which
% lets two defects share one site in different orientations. Each site holds
% one defect, and each defect then picks one of g_X orientations.

There are $N-(i-1)$ free sites and $g_X$ orientations for the $i$-th defect.
Therefore, there are
\begin{align}
    W = \prod_{i=1}^n g_X\,(N-i+1)=g_X^{\,n}\,\frac{N!}{(N-n)!}
\end{align}
ways to place $n$ distinguishable defects in the supercell. At large $N$, the
Stirling approximation applies,
\begin{align}
    x! \approx \sqrt{2\pi x}\bigg(\frac{x}{e}\bigg)^x,
\end{align}
which gives
\begin{align}
    \frac{N!}{(N-n)!}\approx \sqrt{\frac{N}{N-n}}\,\frac{N^{N}}{(N-n)^{N-n}}\,e^{-n}.
\end{align}
% CHANGED: the exponential is e^{-n}, from e^{-N}/e^{-(N-n)}, not
% e^{1-n/(g_x N)}; the square root follows from sqrt(2 pi N)/sqrt(2 pi (N-n)).
The interesting limit is a very dilute defect concentration, $n\ll N$. The
middle factor is
\begin{align}
    \frac{N^{N}}{(N-n)^{N-n}}=N^n\bigg(\frac{N}{N-n}\bigg)^{N-n}
    =N^n\bigg(1-\frac{n}{N}\bigg)^{-(N-n)}\approx N^n e^{n},
\end{align}
% CHANGED (added step): shows where the e^{n} comes from, which cancels the
% e^{-n} of the Stirling approximation.
and the square root tends to one, so
\begin{align}
    \lim_{n/N\rightarrow 0}\, g_X^{\,n}\sqrt{\frac{N}{N-n}}\,\frac{N^{N}}{(N-n)^{N-n}}\,e^{-n}= (g_XN)^n.
\end{align}

If the defects are identical, permuting them leaves the configuration
unchanged. There are $n!$ permutations of one distinct configuration, so the
number of ways to put $n$ indistinguishable defects in the crystal is
\begin{align}
    W =g_X^{\,n}\,\frac{N!}{(N-n)!\,n!}\approx \frac{(g_XN)^n}{n!}.
\end{align}

\paragraph{Configurational entropy.}
The \textcolor{blue}{configurational entropy} of the defect system at fixed
internal energy, crystal volume and particle number is
\begin{align}
    S_{\text{conf}}&= k_B\ln W.
\end{align}
Using the Stirling formula
\begin{align}
    \ln n! \approx n\ln n - n
\end{align}
% CHANGED: the draft wrote n! = n ln n - n; the formula is for ln n!.
gives
\begin{align}
    S_{\text{conf}} &= k_B \big[n\ln (g_XN)-n\ln n +n\big].
\end{align}
As a function of the defect concentration $c=n/N$,
\begin{align}
    S_{\text{conf}} &= k_B \big[cN\ln (g_XN)-cN\ln (cN) +cN\big] \\
    &= Nk_B \big[c\ln g_X - c \ln c + c\big].
\end{align}
The entropy is proportional to $N$, so the configurational entropy per site
is
% CHANGED: the draft justified the division by N with "each defect site is
% independent so that entropy is additive"; dividing by N only expresses the
% result per site, because S_conf is extensive.
\begin{align}
    s_{\text{conf}}(c) &= k_B \big[c\ln g_X - c \ln c + c\big].
\end{align}
Without the dilute limit, the same steps applied to the exact $W$ give
$s_{\text{conf}}(c)=-k_B[c\ln c+(1-c)\ln(1-c)]+k_Bc\ln g_X$, which reduces
to the expression above for $c\ll1$.

\paragraph{Restoring the temperature.}
The DFT calculation of Section~1 neglects the temperature dependence of the
defect formation energy by setting the temperature to $0$~K. This section
restores the temperature, with the assumption that the energy difference
$\Delta U$ between the defective and the pristine crystal does not depend on
the temperature.

At finite temperature, the free energy contains the change in entropy of the
defect and of the bulk. There are three contributions, the electronic, the
vibrational and the configurational entropy. For now, only the
configurational entropy is considered.

\paragraph{The ensemble.}
Section~1 places one defect at a fixed site and exchanges atoms and electrons
with reservoirs at fixed $(T,V,\{\mu_i\},E_F)$, so the formation energy
$E^f[X^q]=\Phi[X^q]-\Phi_{\text{bulk}}$ is a change of the grand potential.
The same ensemble applies to a crystal with $n$ defects.
% CHANGED: the draft switched to a Gibbs energy G[bulk], G[X^q] and a G^f.
% The formation energy stays E^f and the potential stays Phi, as in
% Sections 1 and 2. The fixed-pressure question is in the last paragraph.
The defects are far apart and do not interact, so every arrangement has the
grand potential $\Phi_{\text{bulk}}+nE^f[X^q]$, and every arrangement is
equally probable. The sum over the $W$ arrangements in the grand partition
function gives
\begin{align}
    \Phi(n)=\Phi_{\text{bulk}}+n\,E^f[X^q]-k_BT\ln W
    =\Phi_{\text{bulk}}+n\,E^f[X^q]-T\,S_{\text{conf}}.
\end{align}
% CHANGED: replaces n G^f = U[X^q] - U[bulk] + T S[conf]. The configurational
% entropy belongs to the arrangement of all defects, not to one defect, so
% the formation energy of one defect does not contain S_conf, and the sign of
% the T S_conf term is negative.
\textcolor{red}{Here, $\varphi\equiv\Phi/N$ denotes the grand potential per
site.} Per site,
\begin{equation}
    \boxed{\varphi(c)=\varphi_{\text{bulk}}+c\,E^f[X^q]-T\,s_{\text{conf}}(c).}
\end{equation}

\paragraph{The equilibrium concentration.}
The defects form and disappear freely, so $c$ is an internal variable, and
$\varphi$ is a minimum at equilibrium. With the dilute $s_{\text{conf}}$,
\begin{align}
    \frac{\partial\varphi}{\partial c}
    =E^f[X^q]-T\,\frac{\partial s_{\text{conf}}}{\partial c}
    =E^f[X^q]-k_BT\big[\ln g_X-\ln c\big]=0,
\end{align}
so
\begin{equation}
    \boxed{c_{\text{eq}}=g_X\,e^{-E^f[X^q]/k_BT}.}
\end{equation}

\paragraph{Fixed pressure instead of fixed volume.}
Experiments are done at fixed pressure $P$, not at fixed volume. With $P$ as
the natural variable instead of $V$, the potential gains the term $PV$, and
the formation energy gains $P\,\Delta V$, where $\Delta V$ is the volume
change on forming the defect. At ambient pressure,
$P\,\Delta V\approx10^5~\text{Pa}\times10^{-29}~\text{m}^3\approx10^{-5}$~eV,
negligible against formation energies of the order of 1~eV. The results of
this section therefore hold at fixed pressure as well. Freysoldt \textit{et
al.} write this potential as $G$ and the formation energy at finite
temperature as $G^f$, which correspond to $\Phi$ and $E^f$ in these notes.
